\documentclass[12pt,a4paper]{article}

% Chỉ là mã nguồn LaTeX; chưa biên dịch theo yêu cầu của sinh viên.
% Các khung hình là vị trí để sinh viên tự thay bằng \includegraphics.
% Khi biên dịch sau này: pdfLaTeX và bộ hỗ trợ tiếng Việt T5/babel.
\usepackage[T5]{fontenc}
\usepackage[utf8]{inputenc}
\usepackage{lmodern}
\usepackage[vietnamese]{babel}
\usepackage[a4paper,margin=2.5cm]{geometry}
\usepackage{graphicx}
\graphicspath{{figures/}{images/}}
\usepackage{array,booktabs,float,tabularx,longtable}
\usepackage{amsmath,amssymb}
\usepackage[dvipsnames]{xcolor}
\usepackage{tikz}
\usetikzlibrary{arrows.meta,positioning,shapes.geometric,fit}
\usepackage{subcaption,placeins}
\captionsetup{font=small,labelfont=bf,skip=4pt}
\captionsetup[subfigure]{font=footnotesize,skip=2pt}
\usepackage{xurl}
\usepackage{hyperref}
\linespread{1.25}
\hypersetup{
  colorlinks=true,linkcolor=black,urlcolor=blue,citecolor=black,unicode=true,
  pdftitle={Bài thực hành 07 - Phân vùng bản đồ và điều hướng robot bằng giọng nói},
  pdfauthor={Nguyễn Văn Tổng}
}
\setlength{\parindent}{0pt}
\setlength{\parskip}{6pt}
\setlength{\emergencystretch}{3em}
\renewcommand{\arraystretch}{1.2}
\newcolumntype{Y}{>{\raggedright\arraybackslash}X}
\newcommand{\code}[1]{\texttt{\detokenize{#1}}}
\newcommand{\projectpath}{\path{/home/admin1/hri_ws/week7_voice_navigation_ws}}
\newcommand{\repourl}{https://github.com/VTongg159/hri-turtlebot4-voice-navigation}
\newcommand{\placeholderfigure}[3]{%
  \begin{figure}[H]
    \centering
    \fbox{\begin{minipage}[c][#1][c]{0.91\linewidth}
      \centering\textcolor{gray}{VỊ TRÍ CHÈN HÌNH}\\[6pt]
      \parbox{0.9\linewidth}{\centering\small #2}
    \end{minipage}}
    \caption{#2}\label{#3}
  \end{figure}%
}
\definecolor{flowblue}{RGB}{232,241,250}
\definecolor{flowgreen}{RGB}{232,246,236}
\definecolor{floworange}{RGB}{252,241,224}
\definecolor{linegray}{RGB}{75,75,75}
\tikzset{
  block/.style={draw=linegray,rounded corners,align=center,
    minimum height=10mm,text width=3.2cm,fill=flowblue},
  sensor/.style={block,fill=floworange},
  process/.style={block,fill=flowgreen},
  arrow/.style={-{Latex[length=2.2mm]},thick,draw=linegray}
}

\begin{document}
\pagenumbering{gobble}
\begin{titlepage}
  \centering
  \vspace*{0.8cm}
  {\Large\bfseries ĐẠI HỌC QUỐC GIA HÀ NỘI\par}
  \vspace{0.25cm}
  {\Large\bfseries TRƯỜNG ĐẠI HỌC CÔNG NGHỆ\par}
  \vspace{1.1cm}
  % Thay khung dưới bằng logo nhà trường do sinh viên tự chèn.
  \fbox{\parbox[c][3.5cm][c]{4.5cm}{\centering Vị trí chèn\\logo nhà trường}}
  \vspace{1.1cm}
  {\Huge\bfseries BÀI THỰC HÀNH 07\par}
  \vspace{0.4cm}
  {\LARGE\bfseries Phân vùng bản đồ và điều hướng robot bằng giọng nói\par}
  \vspace{1.3cm}
  \begin{flushleft}
    \renewcommand{\arraystretch}{1.5}
    \begin{tabular}{@{}>{\bfseries\raggedright\arraybackslash}p{5cm}p{9cm}@{}}
      Họ và tên: & Nguyễn Văn Tổng \\
      Mã sinh viên: & 23020766 \\
      Lớp: & K68E-RE \\
      Giảng viên hướng dẫn: & PGS. TS. Hoàng Văn Xiêm \\
                            & KS. Nguyễn Quốc Bảo \\
    \end{tabular}
  \end{flushleft}
  \vfill
  {\large Hà Nội, 2026\par}
\end{titlepage}
\pagenumbering{roman}
\tableofcontents
\newpage
\listoffigures
\listoftables
\newpage
\pagenumbering{arabic}

\section{Giới thiệu và mục tiêu}
Bài thực hành xây dựng hệ thống cho phép người dùng ra lệnh tiếng Việt để
TurtleBot~4 xác định khu vực trên bản đồ và tự di chuyển trong mô phỏng.
Chuỗi xử lý gồm thu âm, nhận dạng tiếng nói (ASR), phân tích câu lệnh,
tra cứu bản đồ ngữ nghĩa và điều hướng bằng ROS~2 Nav2.
Đầu ra của segmentation phải liên kết được với tọa độ điều hướng; việc
gán tên phòng phải cho phép cả gọi theo ID và gọi theo tên chức năng.

Các mục tiêu cụ thể là: đọc bản đồ occupancy dạng PGM/YAML; phân vùng và
lưu polygon, tên, aliases, goal; chuyển đổi chính xác pixel--map; xử lý
lệnh một đích, nguồn--đích và nhiều waypoint; từ chối lệnh sai; kiểm tra
đường đi và theo dõi trạng thái thực tế của robot.
Kiểm thử được tách thành unit test, kiểm thử tích hợp bằng văn bản,
kiểm thử Nav2 trong mô phỏng và kiểm thử end-to-end bằng microphone thật.

Môi trường cuối là PAIRS Robotics Research Lab, kích thước $28\times20$~m,
gồm năm phòng nghiên cứu và một sảnh. Các số liệu trong báo cáo lấy từ
các lần chạy đã lưu; những kết quả của scene campus cũ được ghi rõ phạm
vi và không dùng để khẳng định scene PAIRS đã vượt qua cùng phép thử.

\section{Môi trường mô phỏng và kiến trúc hệ thống}
\subsection{Môi trường và các vùng ngữ nghĩa}
Hệ thống chạy trên Ubuntu~24.04, ROS~2 Jazzy, Gazebo Harmonic và Nav2.
Các model Unitree G1/Go2, Universal Robots UR5e/UR3e, X500 và TurtleBot
được đóng gói cùng tài nguyên hiển thị. Logo PAIRS Lab và ĐHQGHN dùng
ảnh gốc được cung cấp. Nguồn và giấy phép bên thứ ba được lưu riêng.

\begin{table}[H]
\centering\small
\caption{Các vùng ngữ nghĩa và goal trong hệ tọa độ map.}
\label{tab:regions}
\begin{tabularx}{\linewidth}{@{}cYrrr@{}}
\toprule
ID & Khu vực và thiết bị & $x$ (m) & $y$ (m) & yaw (rad) \\
\midrule
A & UAV: X500, lồng thử bay, bãi đáp & $-9.00$ & $3.65$ & $\pi/2$ \\
B & Tay máy: UR5e, UR3e, bàn thao tác & $0.00$ & $3.65$ & $\pi/2$ \\
C & Sảnh PAIRS Lab, logo, lối kết nối & $0.00$ & $0.00$ & $0$ \\
D & HRI/humanoid: Unitree G1 & $9.00$ & $3.65$ & $\pi/2$ \\
E & Robot bốn chân: Unitree Go2 & $-7.00$ & $-3.65$ & $-\pi/2$ \\
F & Robot di động: TurtleBot, dock & $7.00$ & $-3.65$ & $-\pi/2$ \\
\bottomrule
\end{tabularx}
\end{table}

Các robot trưng bày là model tĩnh. Báo cáo không tuyên bố G1/Go2 đang đi,
X500 đang bay hoặc UR đang thao tác. Một TurtleBot~4 mô phỏng riêng là
robot di chuyển bằng Nav2. Vách, cửa, bàn và bục có collision; lối đi
và vị trí goal được bố trí để robot điều hướng tiếp cận các phòng.

\placeholderfigure{6cm}{Toàn cảnh PAIRS Lab và năm khu nghiên cứu.
Gợi ý chèn ảnh pairslab\_overview.png hoặc environment\_overview.png.}{fig:overview}
\placeholderfigure{5cm}{Các góc UAV, tay máy, Unitree G1, Unitree Go2,
TurtleBot và hai logo. Có thể ghép các ảnh chụp Gazebo thật thành nhiều ô.}{fig:rooms}

\subsection{Các module và luồng dữ liệu}
\begin{table}[H]
\centering\small
\caption{Vai trò của các module chính trong package tb4\_voice\_navigation.}
\begin{tabularx}{\linewidth}{@{}lY@{}}
\toprule
Module & Chức năng \\
\midrule
\code{map_segmentation.py} & Hiển thị bản đồ, editor polygon/goal và xuất gợi ý watershed. \\
\code{semantic_map.py} & Đọc occupancy, raster hóa vùng, kiểm tra goal và kết nối. \\
\code{coordinate_utils.py} & Chuyển đổi tâm pixel, tọa độ map và quaternion yaw. \\
\code{voice_recognition_node.py} & Thu ALSA, Whisper tiếng Việt và kiểm soát chất lượng âm thanh. \\
\code{command_parser.py} & Chuẩn hóa tiếng Việt, phân tích cú pháp và khớp aliases. \\
\code{voice_navigation_node.py} & Kiểm tra định vị, lập kế hoạch, gửi và hủy action Nav2 tuần tự. \\
\code{region_visualization_node.py} & Công bố marker vùng và goal cho RViz. \\
\bottomrule
\end{tabularx}
\end{table}

\begin{figure}[H]
\centering
\begin{tikzpicture}[node distance=9mm and 9mm]
\node[sensor] (mic) {Microphone\\ALSA / push-to-talk};
\node[process,right=of mic] (asr) {Whisper small\\transcript tiếng Việt};
\node[block,right=of asr] (parser) {Command parser\\ID và danh sách đích};
\node[block,below=of parser] (semantic) {Semantic map\\polygon, alias, goal};
\node[process,left=of semantic] (nav) {Nav2 action client\\plan, navigate, cancel};
\node[sensor,left=of nav] (robot) {TurtleBot 4\\Gazebo + AMCL / TF};
\draw[arrow] (mic)--(asr);
\draw[arrow] (asr)--(parser);
\draw[arrow] (parser)--(semantic);
\draw[arrow] (semantic)--(nav);
\draw[arrow] (nav)--(robot);
\end{tikzpicture}
\caption{Luồng xử lý từ tiếng nói tới di chuyển. Trạng thái action và TF
được ghi lại để nghiệm thu từng nhiệm vụ.}
\label{fig:architecture}
\end{figure}

Transcript được công bố trên \code{/voice/transcript}; parser tạo dữ
liệu có cấu trúc trên \code{/voice/parsed_command}. Mỗi lệnh có
\code{command_id} để đối chiếu với \code{/voice/navigation_status}.
Các trạng thái ASR xuất hiện trên \code{/voice/asr_status}.
Lệnh nhập bàn phím sử dụng cùng parser/controller nhưng bỏ qua microphone
và ASR, vì vậy kết quả bàn phím không thay thế được kiểm thử T15.

\section{Phương pháp phân vùng bản đồ}
\subsection{Nguồn occupancy và phân loại ô}
Bản đồ \code{pairs_lab.pgm} có $580\times420$ ô, độ phân giải
$r=0.05$~m/ô, origin $(-14.5,-10.5,0)$. Phần bản đồ phủ $29\times21$~m,
bao gồm biên ngoài của khu lab $28\times20$~m.
Đây là bản đồ prior sinh từ cùng geometry collision của world,
\textbf{không phải bản đồ thu được bằng SLAM}. Bản YAML lưu đường dẫn
ảnh, resolution, origin, mode trinary và các ngưỡng occupancy.

Với ảnh xám $I(c,v)\in[0,255]$ và \code{negate=0}, giá trị tương ứng
xác suất chiếm chỗ được tính:
\begin{equation}
 p_{\mathrm{occ}}(c,v)=1-\frac{I(c,v)}{255}.
\end{equation}
Ô được phân thành free nếu $p_{\mathrm{occ}}<0.196$, occupied nếu
$p_{\mathrm{occ}}>0.65$; các giá trị còn lại là unknown.
Unknown và vùng ngoài bản đồ đều không được xem là không gian đi được.

\placeholderfigure{5cm}{Bản đồ occupancy gốc pairs\_lab.pgm, chưa tô màu
ngữ nghĩa; chú thích free, occupied, unknown và vị trí xuất phát.}{fig:originalmap}

\subsection{Phân vùng polygon có giám sát}
Phương pháp dùng trong scene cuối là \textbf{phân vùng polygon và gán
nhãn có giám sát}. Ranh giới A--F được xác định từ bố cục các phòng;
generator sinh các polygon và metadata tương ứng. Công cụ editor cho
phép bổ sung đỉnh bằng chuột, hoàn tác, đổi tên/aliases và điều chỉnh goal.
Không sử dụng mô hình học sâu để tự suy ra chức năng phòng từ ảnh occupancy.

Với polygon vùng $P_i$, mặt nạ vùng hiển thị được định nghĩa:
\begin{equation}
 M_i(c,v)=\mathbf{1}[(c,v)\in P_i]\;\mathbf{1}[(c,v)\in F],
\end{equation}
trong đó $F$ là tập ô free. Việc lấy giao với free giữ nguyên tường và
vật cản nằm trong polygon phòng. Các bước thực hiện là:
\begin{enumerate}
  \item Đọc và kiểm tra metadata, ảnh occupancy, resolution và origin.
  \item Xác định polygon A--F theo kiến trúc; lưu các đỉnh theo chỉ số pixel.
  \item Kiểm tra ít nhất ba đỉnh, tọa độ nguyên hữu hạn, nằm trong ảnh,
        diện tích khác không và không tự giao.
  \item Raster hóa bằng \code{cv2.fillPoly}; tô màu phần polygon giao free.
  \item Gán ID, tên, aliases và một goal hợp lệ; xuất ảnh và YAML.
\end{enumerate}

Ưu điểm của phương pháp là ranh giới và tên phòng có thể kiểm soát,
dễ sửa khi thay bố cục. Hạn chế là cần người thiết kế xác nhận; bản đồ
collision và polygon cùng nguồn thiết kế chưa chứng minh khả năng phân
vùng tự động trên bản đồ SLAM nhiễu của một công trình chưa biết.

\placeholderfigure{6cm}{Ảnh phân vùng A--F: dùng pairslab\_segmentation.png.
Chú thích màu phòng, sảnh, goal, yaw đích và START.}{fig:segmentedmap}

\subsection{Gợi ý watershed tùy chọn}
Module cũng có chế độ \code{--auto} để xuất nhãn gợi ý bằng watershed:
tính distance transform trên free mask; chọn hạt giống tại những điểm có
khoảng cách lớn hơn $\max(2,0.40D_{\max})$ theo đơn vị pixel; mở hình thái
với kernel $3\times3$; tìm connected components làm marker và chạy
\code{cv2.watershed} trên ảnh occupancy chuyển sang ba kênh.
Các ô không free được đặt lại về nhãn nền; kết quả lưu dạng NumPy.

Đây chỉ là công cụ gợi ý cần duyệt thủ công. Scene PAIRS cuối sử dụng
polygon đã xác định theo thiết kế; chưa có tập ground truth độc lập,
IoU hay độ chính xác để kết luận chất lượng watershed. Thuật toán không
tự đặt tên ``UAV'', ``HRI'' hoặc ``Go2'' từ hình học bản đồ.

\subsection{Chọn và kiểm tra goal an toàn}
Goal không được lấy từ centroid polygon một cách mặc định vì centroid
có thể nằm trong vật cản hoặc quá sát tường. Với distance transform $D$
trên free mask đã chặn biên ngoài, clearance bảo thủ theo mét được tính:
\begin{equation}
 d(c,v)=\max\left(0,rD(c,v)-\frac{r}{\sqrt{2}}\right).
\end{equation}
Số hạng trừ bù nửa đường chéo ô lưới. Tập ô ứng viên là free, thuộc vùng,
có $d\geq0.38$~m và nằm trong cùng thành phần liên thông bốn hướng với
anchor $(-11,0)$ trên lưới đã lọc clearance. Nếu cần tự chọn, editor lấy
ô có clearance lớn nhất trong tập ứng viên; goal cấu hình sẵn vẫn phải
vượt qua các kiểm tra này.

Ngưỡng $0.38$~m lớn hơn bán kính robot dùng trong bộ kiểm tra $0.22$~m,
nhưng kiểm tra lưới tĩnh không thay thế kiểm tra footprint, costmap và
động học tại runtime. Trước từng waypoint, Nav2 planner vẫn phải tìm
được một đường đi không rỗng; unknown không được cho phép trong planner.

\section{Semantic mapping và chuyển đổi tọa độ}
\subsection{Cấu trúc metadata và tra cứu tên}
File \code{config/pairs_regions.yaml} lưu frame map, map YAML, anchor,
clearance tối thiểu và danh sách vùng. Mỗi vùng có \code{id},
\code{name}, \code{aliases}, \code{polygon_pixels}, \code{goal}; màu
và vị trí nhãn là trường phục vụ hiển thị. Goal chứa $(x,y,\mathrm{yaw})$
trong frame map. Ví dụ rút gọn:
\begin{verbatim}
map_frame: map
map_yaml: ../maps/pairs_lab.yaml
min_clearance: 0.38
reachable_from_map: [-11.0, 0.0]
regions:
  - id: A
    name: "Phong UAV"  # File that luu ten tieng Viet co dau
    aliases: ["uav", "khu uav", "phong drone"]
    polygon_pixels: [[13,14], [191,14], [191,166], [13,166]]
    goal: {x: -9.0, y: 3.65, yaw: 1.5707963267948966}
\end{verbatim}
Đoạn trên chỉ minh họa một vùng; file bàn giao chứa đủ A--F và aliases
tiếng Việt có dấu. Mỗi alias sau chuẩn hóa phải ánh xạ duy nhất tới một
ID; alias không tồn tại hoặc trùng giữa nhiều vùng đều bị từ chối.

\subsection{Biến đổi tâm pixel sang hệ map}
Gọi $(c,v)$ là cột và hàng tính từ góc trên trái ảnh, $H$ là chiều cao
ảnh, $r$ là resolution và origin là $(x_0,y_0,\theta_0)$.
Tọa độ cục bộ tại \textbf{tâm ô}, có đảo chiều trục ảnh, là:
\begin{equation}
 u=(c+0.5)r,\qquad w=(H-v-0.5)r.
\end{equation}
Tọa độ ROS được tính bằng:
\begin{equation}
 \begin{bmatrix}x\\y\end{bmatrix}=
 \begin{bmatrix}x_0\\y_0\end{bmatrix}+
 \begin{bmatrix}\cos\theta_0&-\sin\theta_0\\
                 \sin\theta_0&\cos\theta_0\end{bmatrix}
 \begin{bmatrix}u\\w\end{bmatrix}.
\end{equation}
Với $\Delta x=x-x_0$, $\Delta y=y-y_0$, phép nghịch đảo là:
\begin{align}
 c&=\frac{\cos\theta_0\Delta x+\sin\theta_0\Delta y}{r}-0.5,\\
 v&=H-\frac{-\sin\theta_0\Delta x+\cos\theta_0\Delta y}{r}-0.5.
\end{align}
Khi truy vấn một ô, chương trình lấy $\lfloor c+0.5\rfloor$ và
$\lfloor v+0.5\rfloor$, sau đó kiểm tra biên. Yaw được chuyển sang
quaternion phẳng $(0,0,\sin(\mathrm{yaw}/2),\cos(\mathrm{yaw}/2))$.
Unit test kiểm tra round-trip với origin có quay; ngưỡng sai số toán học
$10^{-10}$ pixel trong các mẫu thử không phải sai số định vị AMCL.

\placeholderfigure{4.5cm}{Sơ đồ quy ước pixel ở góc trên trái, trục y ảnh
hướng xuống, trục map hướng lên và origin của bản đồ.}{fig:coordinates}

\section{Nhận dạng giọng nói và phân tích câu lệnh}
\subsection{Thu âm và Whisper tiếng Việt}
ASR sử dụng \code{faster-whisper}, mô hình Whisper small đa ngôn ngữ
chạy local trên CPU với \code{compute_type=int8}, bốn CPU threads.
Không dùng TTS hoặc transcript kỳ vọng để thay đầu vào microphone.
Âm thanh được thu bằng ALSA \code{arecord}, mono, PCM signed 16-bit,
$16$~kHz. Push-to-talk cho phép Enter bắt đầu và Enter dừng; một đoạn
tự dừng sau tối đa $12$~s. Hệ thống nhận dạng sau khi kết thúc đoạn thu,
chưa phải ASR streaming liên tục.

\begin{table}[H]
\centering\small
\caption{Cấu hình lọc và giải mã ASR đã triển khai.}
\begin{tabularx}{\linewidth}{@{}lY@{}}
\toprule
Tham số & Giá trị và vai trò \\
\midrule
Thời lượng tối thiểu & $0.4$~s; đoạn ngắn hơn trả NO\_AUDIO. \\
RMS tối thiểu & $0.008$ trên mẫu chuẩn hóa về khoảng $[-1,1]$. \\
Clipping tối đa & Tỷ lệ mẫu có $|a|>0.995$ không vượt $0.08$. \\
Language, beam, temperature & \code{vi}, beam size $5$, temperature $0$. \\
VAD & Bật; min speech $250$~ms, min silence $400$~ms. \\
Điểm segment & avg\_logprob $\geq-0.8$, no\_speech\_prob $\leq0.5$. \\
Giới hạn văn bản & Không rỗng, tối đa $240$ ký tự, ít nhất hai ký tự chữ. \\
\bottomrule
\end{tabularx}
\end{table}

Với $N$ mẫu chuẩn hóa $a_n$, mức năng lượng đo bằng
$\mathrm{RMS}=\sqrt{N^{-1}\sum_{n=1}^{N}a_n^2}$.
Các ngưỡng trên là điều kiện kỹ thuật chấp nhận tín hiệu, không bảo đảm
mọi từ đều được nhận dạng đúng. Nếu không có segment đủ tin cậy, node
phát NO\_SPEECH và không publish transcript để điều hướng. Khi thiết bị
âm thanh không khả dụng, lỗi ALSA được chuyển thành trạng thái lỗi.
WAV và audit log có thể lưu theo cấu hình để đối chiếu từng lượt.

\subsection{Chuẩn hóa và grammar}
Parser chạy độc lập với ASR. Văn bản được chuyển về chữ thường, chuẩn
hóa Unicode NFD, loại dấu kết hợp, chuyển ``đ'' thành ``d'', thay dấu câu
bằng khoảng trắng và gộp khoảng trắng. Chuẩn hóa giúp ``cho robot''
khớp alias ``chó robot'', nhưng không tự sửa từ sai thành một vùng gần nhất.

Grammar nhận lệnh có động từ di chuyển rõ ràng, tách các phần nối
``đến'', ``tới'', ``qua'', ``rồi đến'', ``sau đó đến'' và tra cứu toàn bộ
các đích trước khi gửi nhiệm vụ. Lệnh ``từ'' phải có ít nhất nguồn và
đích. Tối đa $12$ waypoint/lệnh; lệnh dừng được phân tích riêng.

\begin{table}[H]
\centering\small
\caption{Ví dụ ánh xạ câu lệnh trên scene PAIRS Lab.}
\begin{tabularx}{\linewidth}{@{}YY@{}}
\toprule
Câu lệnh & Kết quả parser \\
\midrule
Đi đến khu vực A & single, [A] \\
Đi đến phòng UAV & single, [A] \\
Đi đến phòng chó robot & single, [E] \\
Đi từ A đến B & source\_destination, [A,B] \\
Đi từ phòng UAV qua hành lang đến phòng robot hình người & multi\_waypoint, [A,C,D] \\
Đi đến khu vực Z & UNKNOWN\_REGION; không tạo route hợp lệ \\
Dừng lại & cancel \\
\bottomrule
\end{tabularx}
\end{table}

Tên chung ``phòng thí nghiệm'' trong đề gốc có alias trong cấu hình
campus dùng cho unit test. Trong scene PAIRS có năm lab khác chức năng,
tên chung này bị từ chối; demo semantic dùng tên cụ thể ``phòng UAV''.
Đây là thay đổi có chủ đích của scene cuối, không kết luận nguyên văn
câu T06 đã được chấp nhận ở scene PAIRS.

\placeholderfigure{6cm}{Ảnh chụp trực tiếp terminal trong lượt thử ngày
  10/11/2026: push-to-talk, transcript ASR, parser xác định vùng A cùng
  command\_id và kết quả Nav2 SUCCEEDED của câu lệnh đến phòng UAV.}{fig:asr}

\section{Tích hợp ROS 2 Nav2 và xử lý lỗi}
Controller kiểm tra cấu trúc request và tất cả region trước khi nhận
nhiệm vụ. Trạng thái chính là IDLE/VALIDATING, PLANNING, NAVIGATING,
CANCELING và trạng thái kết thúc SUCCEEDED/FAILED/CANCELED.
Robot chỉ thực hiện một nhiệm vụ tại một thời điểm; lệnh mới khi busy
bị từ chối, không đẩy chồng thêm action.

Trước khi đi, controller yêu cầu AMCL pose trong frame map, TF
\code{map -> base_link} còn mới không quá $2$~s và các phần tử covariance
ở chỉ số $0,7,35$ hữu hạn, không âm, không quá $0.5$ theo cấu hình.
Đây là ngưỡng trên từng phần tử; thành phần vị trí có đơn vị m$^2$,
thành phần góc có đơn vị rad$^2$. Hai action server phải sẵn sàng.

Mỗi waypoint được xử lý theo thứ tự:
\begin{enumerate}
  \item Kiểm tra lại goal và tạo PoseStamped trong frame map.
  \item Gọi \code{/compute_path_to_pose}, yêu cầu terminal success,
        error code bằng $0$ và path không rỗng.
  \item Gọi \code{/navigate_to_pose}; nhận feedback và terminal result thật.
  \item Chỉ tăng chỉ số waypoint khi action status bằng $4$
        (SUCCEEDED) và error code bằng $0$.
\end{enumerate}
Với lệnh nguồn--đích, đường thực hiện là current$\to$source$\to$destination;
chỉ bỏ qua source nếu robot đang cách goal nguồn không quá $0.30$~m.
Điều này tránh hiểu sai ``đi từ A đến B'' thành đi thẳng tới B.

Server timeout là $10$~s, planning timeout $15$~s và goal timeout $180$~s.
Khi người dùng hủy, hết giờ hoặc mất định vị, controller yêu cầu cancel
và giữ busy cho tới terminal result. Nếu hủy chưa được xác nhận sau
$10$~s, hệ thống báo tình trạng và tiếp tục giữ khóa; không giả định
robot đã dừng chỉ vì đã gửi cancel request.
Service \code{/voice/cancel} cho phép hủy mà không phải chờ ASR của một
câu mới; hủy bằng giọng nói có độ trễ thu và giải mã đoạn âm thanh.

\section{Kết quả thử nghiệm}
\subsection{Phạm vi, tiêu chí và bằng chứng}
Kết quả dưới đây phân biệt ba nguồn: unit/integration test của package;
các lần Nav2 chạy trên PAIRS Lab; các lượt microphone thật.
Bằng chứng gồm XML test, goal validation, JSON kết quả và event log có
command\_id. Terminal SUCCEEDED được đối chiếu với action status/error
code, không dựa riêng vào việc robot trông có vẻ đã đến gần phòng.

\begin{table}[H]
\centering\small
\caption{Các lớp kiểm chứng và phạm vi kết luận.}
\begin{tabularx}{\linewidth}{@{}lYY@{}}
\toprule
Lớp thử & Kết quả đã ghi & Giới hạn \\
\midrule
Unit test & 87 test, không lỗi/thất bại/bỏ qua & Gồm fixture campus và test riêng PAIRS; không phải 87 lượt robot chạy. \\
Goal PAIRS & Sáu goal hợp lệ, clearance $\geq0.38$~m & Kiểm tra geometry/lưới tĩnh; chưa bao quát mọi vật cản động. \\
Nav2 PAIRS & Sáu đích A,B,D,C,E,F đạt & Một lượt kiểm chứng; không suy ra độ tin cậy dài hạn. \\
Parser/controller PAIRS & Lệnh văn bản tới F đạt & Không dùng microphone, không tính là T15. \\
Microphone/T15 & Ba câu có lượt đạt sau tổng sáu lần ghi & Có retries; chưa đo WER hay nhiều người nói. \\
Video bổ sung & Một câu voice tới A và tour đủ năm phòng & Có cắt dựng và đoạn nối sau mất định vị. \\
\bottomrule
\end{tabularx}
\end{table}

Kiểm tra scan alignment dùng $10$ scan, $3600$ endpoint, TF tại timestamp
scan. Trung vị khoảng cách endpoint tới ô occupied gần nhất là $0$~m,
p95 là $0.1006$~m, p99 là $0.8515$~m; không có endpoint ngoài bản đồ.
Tiêu chí đặt trước là median $\leq0.15$~m và p95 $\leq0.40$~m nên lượt
đo đạt. Đây là kiểm tra độ phù hợp scan/map trong mô phỏng, không phải
đánh giá SLAM hoặc độ chính xác segmentation bằng ground truth độc lập.

\subsection{Goal và điều hướng thực trong PAIRS Lab}
\begin{table}[H]
\centering\small
\caption{Lượt Nav2 kiểm chứng sáu đích; nguồn navigation.json của PAIRS.}
\label{tab:navresults}
\begin{tabular}{@{}crrrcl@{}}
\toprule
Đích & Clearance (m) & Thời gian (s) & $e_{fb}$ (m) & Recovery & Kết quả \\
\midrule
A & 1.015 & 12.06 & 0.223 & 0 & SUCCEEDED \\
B & 1.837 & 41.00 & 0.259 & 0 & SUCCEEDED \\
D & 1.837 & 43.43 & 0.251 & 0 & SUCCEEDED \\
C & 1.865 & 43.35 & 0.025 & 0 & SUCCEEDED \\
E & 1.700 & 39.06 & 0.248 & 0 & SUCCEEDED \\
F & 1.878 & 45.05 & 0.241 & 0 & SUCCEEDED \\
\bottomrule
\end{tabular}
\end{table}
Thời gian là wall time của từng đoạn trong thứ tự A$\to$B$\to$D$\to$C$\to$E$\to$F$.
$e_{fb}$ được tính từ \emph{feedback vị trí cuối được lưu} tới goal;
feedback có chu kỳ nên đây không phải phép đo pose chính xác tại thời
điểm action kết thúc. Các giá trị gần hoặc nhỉnh hơn $0.25$~m không được
dùng để suy ngược goal checker thất bại. Cột kết quả dựa trên terminal
result của Nav2. Lượt này không có recovery; điều đó không loại trừ lỗi
định vị đã gặp ở phiên quay khác.

% Hinh 13: submission/figures/figure13_rviz_UAV_succeeded.png
  \placeholderfigure{6cm}{Ảnh chụp RViz sau khi robot đến phòng UAV (A):
  đường lập kế hoạch Nav2, marker đích, TF robot và trạng thái SUCCEEDED.
  Đường xanh là đường lập kế hoạch được giữ lại, không phải quỹ đạo đã đo.}{fig:navigation}

\subsection{Ba câu tiếng Việt qua microphone thật}
Ngày 09/11/2026, người dùng thực hiện ba câu trong một phiên có lưu WAV
gốc, transcript, parser và Nav2. Kết quả lượt đạt được giữ nguyên như
ASR xuất ra, không chỉnh sửa thành câu mong đợi:

\begin{table}[H]
\centering\small
\caption{Kết quả ba câu tiếng Việt trong phiên nghiệm thu T15.}
\label{tab:voice}
\begin{tabularx}{\linewidth}{@{}cYcc@{}}
\toprule
Câu & Transcript thực tế của lượt đạt & Route & Lượt đạt \\
\midrule
1 & Đi đến phòng UAV & A & 1 \\
2 & Đi đến phòng cho robot & E & 3 \\
3 & Đi từ phòng UAV qua Hành Lang đến phòng robot Hình Người & A$\to$C$\to$D & 2 \\
\bottomrule
\end{tabularx}
\end{table}

Tổng cộng có sáu lượt ghi, ba lượt thành công. Hai lượt trả transcript
sai: ``đi đến phòng chó rối'' và ``đi từ phẩm UAV qua hành lang đến phẩm
robot hình người''; parser trả UNKNOWN\_REGION và không gửi goal mới.
Một lượt trả NO\_SPEECH, không phát transcript. Lượt đạt câu~2 thiếu dấu
ở ``cho robot'' nhưng khớp alias sau chuẩn hóa dấu.

\begin{table}[H]
\centering\small
\caption{Thời gian xử lý của các lượt voice đạt, đo theo timestamp event.}
\begin{tabular}{@{}crrl@{}}
\toprule
Câu & ASR phase (s) & Transcript$\to$mission end (s) & Nav2 \\
\midrule
1 & 1.800 & 56.304 & SUCCEEDED \\
2 & 1.600 & 22.496 & SUCCEEDED \\
3 & 1.900 & 78.059 & SUCCEEDED \\
\bottomrule
\end{tabular}
\end{table}
ASR phase là khoảng TRANSCRIBING$\to$TRANSCRIPT đã publish, không bao
gồm toàn bộ thời gian người nói. Tỷ lệ hoàn thành theo lượt ghi ở phiên
này là $3/6$; chỉ một trong ba câu đạt ngay lượt đầu. Sau thử lại, cả ba
câu có lượt hoàn thành. Các tỷ lệ này không phải word error rate hoặc
độ chính xác ASR tổng quát vì chỉ có một người nói và rất ít mẫu.

\subsection{Ảnh terminal trực tiếp bổ sung ngày 10/11/2026}
Để chụp Hình~\ref{fig:asr}, thực hiện một phiên microphone mới với bộ
quay màn hình bật trước khi thu âm. Lần đầu ASR nhận ``Đi đến phòng UAB'',
parser trả UNKNOWN\_REGION và không gửi goal. Lần sau ASR nhận
``đi đến phòng UAV'', parser xác định [A], Nav2 kết thúc SUCCEEDED với
status~4/error~0. Audio lượt đạt dài $4.734$~s; thời gian từ transcript
đến kết thúc nhiệm vụ là $13.519$~s. Các số này thuộc phiên bổ sung,
không thay cho số liệu ba câu ngày 09/11/2026 ở bảng trước.

Ảnh chụp nguyên cửa sổ terminal khi ứng dụng chạy, giữ cả lượt chưa đạt;
video terminal liên tục dài $61.2$~s, $5$~fps, không biên soạn lại nội dung.
Bằng chứng nằm trong \code{submission/evidence/figure7_live_acceptance.json},
log gốc và event ROS cùng command\_id. WAV lưu local, kèm hash đối chiếu.
Sau nhiệm vụ, TF đo được $(-9.067,3.497)$~m nằm trong vùng A; không coi
đây là pose đúng thời điểm terminal action.

\subsection{Video và tuyến đi đủ năm phòng}
Video bàn giao dài khoảng \textbf{4 phút 26 giây}, $1920\times1080$,
$15$~fps, audio AAC. Đã kiểm tra giải mã toàn bộ file không có lỗi.
Các phần gồm nhận diện logo, toàn cảnh, góc cận từng phòng, một câu voice
thật tới UAV và tour Nav2 dùng văn bản tới các phòng còn lại.

Đoạn voice quay bổ sung sử dụng chính microphone của người dùng:
``Đi đến phòng UAV'' $\to$ ASR $\to$ parser A $\to$ Nav2 SUCCEEDED.
Phiên quay mới có hai lần ghi: lần đầu NO\_SPEECH, lần sau đạt. Bản dựng
chỉ giữ lượt đạt; cả hai lượt vẫn có trong raw video/WAV/log local.
Đoạn voice giữ tốc độ thật. Audio cùng lượt được ghép theo timestamp
RECORDING và thời điểm yêu cầu screen capture; có sai số khởi động
thiết bị/capture, không khẳng định đồng bộ chính xác từng frame.

Tour thực tế đạt theo thứ tự A$\to$B$\to$D$\to$C$\to$F$, sau đó E$\to$C$.
Lần đi F$\to$E đầu bị LOCALIZATION\_LOST khi covariance vị trí vượt
ngưỡng $0.5$; controller hủy nhiệm vụ. Sau cập nhật laser lúc đứng yên
bằng \code{/request_nomotion_update}, AMCL hội tụ lại và lần đi tiếp
E$\to$C đạt. Không teleport, không hạ ngưỡng định vị. Bản dựng bỏ đoạn
thử E thất bại và chờ khắc phục, ghi rõ nhãn đoạn nối sau AMCL; raw log
và clip vẫn giữ. Do đó đây là \textbf{video có cắt dựng}, không phải một
lượt liên tục không lỗi.

Phần tour văn bản có nhãn nguồn và tăng tốc $2\times$.
B/C/D/E/F có pose TF dạng số tại terminal success và kiểm tra trong
polygon trong coverage receipt. Vị trí A đã được kiểm tra khi bắt đầu
tour, nhưng receipt quay hiện tại chưa ghi riêng số đo A; video voice
có HUD pose thật. Báo cáo không bổ sung ngược một số đo chưa lưu.

\begin{table}[H]
\centering\small
\caption{Các mốc chính trong video bàn giao.}
\begin{tabularx}{\linewidth}{@{}lY@{}}
\toprule
Mốc & Nội dung \\
\midrule
00:00 & Nhận diện PAIRS/ĐHQGHN, toàn cảnh và năm khu nghiên cứu. \\
01:07 & Voice thật: ``Đi đến phòng UAV'', parser và robot tới A. \\
02:25 & Bắt đầu tour văn bản $2\times$: B, D, C, F. \\
03:53 & Đoạn nối sau cập nhật AMCL, có ghi nhãn cắt dựng. \\
03:54 & Phòng chó robot Go2, đích E. \\
04:03 & Về sảnh C và kết thúc. \\
\bottomrule
\end{tabularx}
\end{table}

% Hinh 14: submission/figures/figure14_voice_video_UAV_succeeded.png
  \placeholderfigure{5cm}{Minh họa từ khung video tại 02:19 ngày
  09/11/2026: robot đến phòng UAV, transcript ASR, trạng thái SUCCEEDED
  và vị trí robot. Cửa sổ Gazebo giữ nguyên màu gốc; bảng thông tin phụ
  được trình bày lại bằng trắng--đen.}{fig:videovoice}
\placeholderfigure{5cm}{Các ảnh đến B/D/F/E và ảnh kết thúc ở sảnh;
chụp sau SUCCEEDED. Không dùng model trưng bày thay robot điều hướng.}{fig:arrivals}

\textbf{Xác nhận phạm vi video:} có một đoạn thể hiện đầy đủ nhận lệnh
giọng nói, xác định đích và robot di chuyển đến đích. Không phải toàn bộ
tour năm phòng đều được gọi bằng microphone. Video hiện tại cũng chưa
chứa riêng tất cả sáu dạng demo trong đề gốc: câu gọi ID, câu nguồn--đích,
lệnh Z sai và multi-waypoint voice chưa được quay đầy đủ trong cùng bộ
clip; trạng thái này được ghi trong kịch bản quay. Phiên ba câu trước
có audio/log cho multi-waypoint nhưng chưa có video cùng lượt.

\subsection{Đối chiếu các test case bắt buộc}
\small
\begin{longtable}{@{}p{0.8cm}p{4.1cm}p{9.2cm}@{}}
\caption{Đối chiếu T01--T15 và phạm vi bằng chứng.}\label{tab:testmatrix}\\
\toprule ID & Nội dung & Kết quả và phạm vi \\
\midrule\endfirsthead
\toprule ID & Nội dung & Kết quả và phạm vi \\
\midrule\endhead
\bottomrule\endfoot
T01 & Đọc map YAML/PGM & PASS unit; đọc cả fixture campus và map PAIRS. \\
T02 & Pixel--map & PASS unit; round-trip gồm origin quay, kiểm tra tâm ô. \\
T03 & Goal trong occupied & PASS unit từ chối; bộ goal PAIRS cũng được validate. \\
T04 & Goal ngoài polygon & PASS unit từ chối, không lấy goal bất kỳ theo tên phòng. \\
T05 & Đi đến khu vực A & PASS parser unit trên fixture; chưa có clip voice riêng nguyên câu này ở scene cuối. \\
T06 & Đi đến phòng thí nghiệm & PASS alias unit với campus. Scene PAIRS từ chối tên chung; semantic voice dùng ``phòng UAV''. \\
T07 & Đi từ A đến B & PASS parser và controller integration lịch sử; chưa chạy lại riêng câu nguồn--đích hai điểm trong PAIRS. \\
T08 & A qua C đến B & PASS grammar unit/integration lịch sử; PAIRS voice đã chạy route A--C--D, không coi là đã quay nguyên route A--C--B. \\
T09 & Khu vực Z & PASS unit từ chối toàn bộ lệnh; chưa quay microphone riêng câu Z. \\
T10 & Thiếu microphone & PASS integration lịch sử bằng thiết bị ALSA không tồn tại: báo ASR\_ERROR, không crash. \\
T11 & Nav2 không sẵn sàng & PASS integration lịch sử: FAILED/NAV2\_NOT\_READY. Chưa lặp lại trên scene PAIRS. \\
T12 & Goal unreachable & PASS simulation lịch sử: vật cản Gazebo thật, status~6/error~208 và không chạy waypoint sau. Chưa lặp lại trong PAIRS. \\
T13 & Hủy & PASS integration lịch sử với action terminal thực. Phiên quay PAIRS cũng ghi nhận cancel do mất định vị; đây không thay thế thử voice cancel riêng. \\
T14 & Lệnh mới khi busy & PASS integration lịch sử: lệnh mới bị từ chối; chưa lặp lại riêng trên PAIRS. \\
T15 & Voice end-to-end & PASS PAIRS: ba câu có lượt đạt sau sáu lần ghi; video mới bổ sung một câu UAV đạt sau hai lần ghi. \\
\end{longtable}
\normalsize
Các test lịch sử được lưu riêng với cấu hình campus. Việc formatter và
unit suite chạy lại không biến các test tích hợp lịch sử thành phép đo
mới trên scene PAIRS; các trường hợp chưa lặp lại vẫn giữ đúng trạng thái.

\section{Sản phẩm bàn giao và tái lập}
\subsection{Mã nguồn, bản đồ và video}
Repository public: \href{\repourl}{VTongg159/hri-turtlebot4-voice-navigation}.
Phần mềm tự viết và tài nguyên bên thứ ba có thông tin giấy phép riêng.
Bộ bàn giao tập trung ở \code{submission/}:
\begin{itemize}
  \item \code{maps/pairs_lab.pgm} và \code{maps/pairs_lab.yaml}: occupancy gốc.
  \item \code{maps/pairs_lab_original.png}: ảnh xem nhanh bản đồ gốc.
  \item \code{maps/pairslab_segmentation.png}: ảnh phân vùng A--F.
  \item \code{config/pairs_regions.yaml}: polygon, tên, aliases và goal;
        cấu trúc thư mục giữ đúng đường dẫn map tương đối.
  \item \code{evidence/}: goal validation, navigation, voice, coverage,
        kết quả test/clean code và manifest hash các sản phẩm công khai.
  \item \code{PAIRS_Lab_Demo_1080p.mp4}: video bản dựng, kèm mốc thời gian.
  \item File LaTeX hiện tại và kịch bản quay nằm ở \code{docs/}; ảnh trong
        báo cáo này là placeholder để sinh viên tự chèn.
\end{itemize}
Audio WAV thô, cache model, môi trường Python, build/install/log và các
dump tiến trình cá nhân không thuộc bộ source public. Video đã xuất có
tiếng nói trong lượt demo được bàn giao theo yêu cầu.

\subsection{Các bước chạy lại}
Sau khi cài ROS~2 Jazzy và các dependency được ghi trong README, tạo
virtual environment kế thừa gói ROS hệ thống; tải model Whisper lần đầu:
\begin{verbatim}
git clone \
  https://github.com/VTongg159/hri-turtlebot4-voice-navigation.git
cd hri-turtlebot4-voice-navigation
source /opt/ros/jazzy/setup.bash
python3 -m venv --system-site-packages .venv
source .venv/bin/activate
python -m pip install -r requirements-asr.txt
python scripts/download_model.py --model small
python -m colcon build --base-paths src --symlink-install
source install/setup.bash
export ROS_DOMAIN_ID=19
ros2 launch tb4_voice_navigation voice_navigation.launch.py start_sim:=true
\end{verbatim}
Chờ MODEL\_READY, AMCL/TF và Nav2 sẵn sàng. Ở terminal đã source cùng
môi trường và domain, chạy \code{ros2 run tb4_voice_navigation push_to_talk},
Enter để nói rồi Enter dừng; đợi nhiệm vụ kết thúc trước câu tiếp theo.
Pose khởi đầu được cấu hình ở $(-11,0,0)$ cho lần spawn mới; không seed
lại pose spawn sau khi robot đã di chuyển sang vị trí khác.

Tái xuất ảnh segmentation bằng lệnh sau; không thay đổi geometry world:
\begin{verbatim}
ros2 run tb4_voice_navigation map_segmentation \
  src/tb4_voice_navigation/config/pairs_regions.yaml \
  --output segmented_map.png
\end{verbatim}
Lệnh \code{--edit} chỉnh polygon/goal; \code{--auto --output suggestions.npy}
xuất gợi ý watershed để người dùng duyệt. Sau khi sửa bản đồ hoặc metadata
cần kiểm tra lại goal và chạy Nav2; receipt cũ không xác nhận geometry mới.

Kiểm tra source sau dọn mã gồm formatter, lint, build và unit suite:
\begin{verbatim}
python -m pip install -r requirements-dev.txt
ruff check src/tb4_voice_navigation scripts
python -m colcon test --packages-select tb4_voice_navigation
python -m colcon test-result --verbose
\end{verbatim}

\section{Hạn chế, hướng phát triển và kết luận}
Segmentation hiện dựa vào thiết kế và duyệt thủ công; chưa có đánh giá
IoU của watershed trên bản đồ SLAM. Cần xây dựng ground truth độc lập,
đánh giá khả năng tách phòng qua cửa hẹp và xử lý nhiễu/unknown.
Goal an toàn trên map tĩnh vẫn cần planner/costmap runtime, đặc biệt khi
thêm người hoặc vật cản động.

ASR mới được đo với một người nói, có lỗi từ vựng và NO\_SPEECH. Cần thử
nhiều giọng, khoảng cách microphone, mức nhiễu; đo WER và tỷ lệ hoàn thành
nhiệm vụ với số mẫu lớn hơn. Grammar/aliases được giới hạn có chủ đích;
tên phòng mơ hồ cần yêu cầu làm rõ thay vì đoán đích. Có thể bổ sung ASR
streaming và cơ chế dừng có độ trễ thấp sau khi kiểm chứng an toàn.

Phiên quay cho thấy AMCL có thể mất độ tin cậy trên đoạn di chuyển dài.
Cần đánh giá lại scan/map, cấu hình AMCL và hình học phân biệt trong
hành lang; lặp nhiều lượt trước khi kết luận độ ổn định. Các model
trưng bày hiện chưa có controller động học/động lực học riêng.

Kết quả đã thực hiện chứng minh luồng microphone thật đến Nav2 và
TurtleBot~4 tới các đích trong PAIRS Lab, cùng khả năng tra cứu semantic
và chạy nhiều waypoint. Bộ sản phẩm gồm source, occupancy, vùng ngữ nghĩa,
validation, video và tài liệu. Các lần chưa đạt và phạm vi test lịch sử
được giữ rõ để việc đánh giá có thể truy vết và tái lập.

\begin{thebibliography}{9}
\bibitem{source}
Nguyễn Văn Tổng, \emph{Bài thực hành 07: PAIRS Lab voice navigation},
mã nguồn và cấu hình triển khai, 2026.
\href{\repourl}{Repository GitHub public}.
\bibitem{evidence}
\emph{Kết quả thực thi PAIRS Lab ngày 09/11/2026}, các file navigation,
goal validation, scan alignment, voice acceptance và video acceptance;
xem thư mục \code{submission/evidence/} trong repository.
\end{thebibliography}
\end{document}
